%
% 33-residuum.tex
%
% (c) 2026 Prof Dr Andreas Müller
%
\subsection{Das Residuum}
Wir betrachten die Laurent-Reihe
\[
f(z)
=
\sum_{k=-n}^\infty a_k(z-z_0)^k
\]
im Punkt $z_0$.
Sei $\gamma$ ein geschlossener Weg in einem Ringgebiet um $z_0$,
 der den Punkt $z_0$ einmal umläuft.
Für die Berechnung des Wegintegrals kann dieser Weg in den kreisförmigen
Weg $\gamma_0(t) = z_0+re^{2\pi it}$ deformiert werden, der den gleichen
Wert für das Integral ergibt.
\begin{align*}
\frac{1}{2\pi i}
\oint_\gamma f(z)\,dz
&=
\frac{1}{2\pi i}
\oint_{\gamma_0} f(z)\,dz
\\
&=
\frac{1}{2\pi i}
\int_0^1
f(z_0+re^{2\pi it}) \dot{\gamma}_0(t)\,dt
=
\frac{1}{2\pi i}
\int_0^1 
f(z_0+re^{2\pi it}) 2\pi ir e^{2\pi it}\,dt
\\
&=
\sum_{k=-n}^\infty \int_0^1 a_k r^k e^{2\pi i k t} re^{2\pi it}\,dt
=
\sum_{k=-n}^\infty \int_0^1 a_k r^{k+1} e^{2\pi i (k+1) t}\,dt
\intertext{Im Fall $k+1= 0$ wird der Integrand konstant, daher behandeln 
wird diesen Fall separat:}
&=
\sum_{\scriptstyle k=-n\atop \scriptstyle k\ne -1}^\infty
a_k r^{k+1} \int_0^1 e^{2\pi i (k+1) t} \,dt
+
a_{-1} \int_0^1 \,dt
\\
&=
\sum_{\scriptstyle k=-n\atop \scriptstyle k\ne -1}^\infty a_k r^{k+1}
\underbrace{\biggl[
\frac{1}{2\pi i(k+1)} e^{2\pi i (k+1) t}
\biggr]_0^1}_{\displaystyle = 0}
\mathstrut
+
a_{-1}
=
a_{-1}.
\end{align*}
Die Integrale für $k\ne -1$ verschwinden, weil der Integrand periodisch
ist und das Integral über eine Periode erstreckt wird.

\begin{definition}[Residuum]
\index{Residuum}%
\index{resz0f@$\operatorname{res}(z_0;f)$}%
Der Koeffizient $a_{-1}$ heisst das \emph{Residuum}
\[
\operatorname{res}(z_0;f)
=
a_{-1}
\]
der Funktion $f(z)$ an der Stelle $z_0$.
\end{definition}

Das Residuum fängt also das Verhalten von Polen erster Ordnung an
der Stelle $z_0$ ein.
Dies wird verdeutlich durch den folgenden Satz.

\begin{satz}
\label{buch:analytisch:laurent:residuum:satz}
Ist $f(z)$ in einer Umgebung von $\alpha$ holomorph und hat $g(z)$
eine Nullstelle erster Ordnung in $\alpha$, dann ist das Residuum
\[
\operatorname{res}\biggl(\alpha;\frac{f(z)}{g(z)}\biggr)
=
\frac{f(\alpha)}{g'(\alpha)}.
\]
\end{satz}

\begin{proof}
Da $g(z)$ eine Nullstelle erster Ordnung an der Stelle $\alpha$
hat, kann man
\[
g(z) = (z-\alpha)\bigl(g'(\alpha)+ r(z)(z-\alpha)\bigr)
\]
schreiben, wobei $r(z)$ eine in einer Umgebung von $\alpha$ holomorphe
Funktion ist, die keine Nullstelle $z=\alpha$ hat.
Damit wird
\[
\frac{f(z)}{g(z)}
=
\frac{f(z)}{(z-\alpha)\bigl(g'(\alpha)+r(z)(z-\alpha)\bigr)}
=
\frac{1}{z-\alpha}
\cdot
\frac{f(z)}{g'(\alpha)+r(z)(z-\alpha)}
=
\frac{1}{z-\alpha}
\cdot
h(z)
\]
mit der in einer Umgebung von $\alpha$ holomorphen Funktion
\[
h(z) = \frac{f(z)}{g'(\alpha)+r(z)(z-\alpha)}
\]
schreiben.
Nach Satz~\ref{buch:analytisch:laurent:satz:laurentreihe} ist
das Residuum das Integral
\begin{align*}
\frac{1}{2\pi i}
\oint_\gamma
\frac{f(z)}{g(z)}
\,dz
&=
\frac{1}{2\pi i}
\oint_\gamma
\frac{1}{z-\alpha}
\frac{f(z)}{g'(\alpha)+r(z)(z-\alpha)}
\,dz
\\
&=
\frac{1}{2\pi i}
\oint_\gamma
\frac1{z-\alpha}
\cdot
h(z)
\,dz.
\intertext{Nach dem Integralsatz~\ref{buch:integration:satz:residuensatz}
ergibt das Integral}
&=
h(\alpha)
=
\frac{f(\alpha)}{g'(\alpha) + r(\alpha)(\alpha - \alpha)}
=
\frac{f(\alpha)}{g'(\alpha)}.
\end{align*}
Dies beweist den Satz.
\end{proof}

\begin{beispiel}
Um das Residuum der Funktion $(1+z)/(1-z)$ an der Stelle $1$ zu
berechnen, kann man den Satz auf die Situation $f(z)=1+z$ und
$g(z)=1-z$ anwenden, und erhält
\[
\operatorname{res}\biggl(1; \frac{1+z}{1-z}\biggr)
=
\frac{f(1)}{g'(1)}
=
\frac{2}{1}
=
2.
\qedhere
\]
\end{beispiel}

Falls der Weg $\gamma$ den Punkt $z_0$ mehrmals umläuft, kann man
das Wegintegral
\begin{equation}
\frac{1}{2\pi i}
\oint_\gamma f(z)\,dz
=
\operatorname{ind}_\gamma(z_0) \cdot \operatorname{res}(z_0;f)
\label{buch:analytisch:laurentreihen:residuum:eqn:integralres}
\end{equation}
aus der Umlaufzahl von $\gamma$ um den Punkt und dem Residuum der
Funktion im Punkt $z_0$ berechnen.

%
% Residuen und Integralberechnung
%
\subsubsection{Residuen und Integralberechnung}
Der Begriff des Residuums ermöglicht manchmal, die Berechnung reeller 
Integrale im Abschnitt~\ref{buch:integration:section:residuen}
noch etwas eleganter zu formulieren.

\begin{beispiel}
\label{buch:analytisch:laurentreihen:bsp:cauchyverteilung}
Bei der Berechnung des Integrals der Funktion
\[
\varphi(z)
=
\frac{1}{z^2+1}
\]
entlang der reellen Achse in
Abschnitt~\ref{buch:integration:residuen:section:beispiel} 
erfolgt mithilfe eines Weges, die die Polstelle $i$ des Integranden
einmal umrundet.
Der Integrand kann als
\[
\varphi(z)
=
\frac{1/(z+i)}{z-i}
=
\frac{f(z)}{g(z)}
\qquad\text{mit}\qquad
\left\{
{\renewcommand{\arraycolsep}{2pt}
\begin{array}{rcl}
f(z) &=& \displaystyle \frac{1}{z+i} \mathstrut
\\[8pt]
g(z) &=& z-i \mathstrut
\end{array}}
\right.
\]
geschrieben werden.
Das Residuum ist nach Satz~\ref{buch:analytisch:laurent:residuum:satz}
daher
\[
\operatorname{res}(i, \varphi)
=
\frac{f(i)}{g'(i)}
=
\frac{1/2i}{1}
=
\frac{1}{2i}.
\]
Für den geschlossenen Pfad $\gamma$ wird nach
\eqref{buch:analytisch:laurentreihen:residuum:eqn:integralres}
\[
\oint_{\gamma} \varphi(z)\,dz
=
2\pi i \operatorname{res}(i,\varphi)
=
2\pi i\cdot \frac{1}{2i}
=
\pi.
\]
Da der Kreisbogen für $R\to\infty$ keinen Beitrag leistet, ist dies auch
der Wert des reellen Integrals.
\end{beispiel}

